\documentclass[blue]{beamer}



\beamertemplatenavigationsymbolsempty
\usetheme{}
% Specify overlay option (for hidden text):
% - invisible
% - transparent
\setbeamercovered{invisible}
\usepackage{xcolor}
  \definecolor{colorLink}{rgb}{0,0.2,0.6}
\usepackage{graphicx}
\usepackage{hyperref}
  \hypersetup{colorlinks,breaklinks,urlcolor=colorLink,linkcolor=colorLink}



\begin{document}



\title{Personal cybersecurity}
\subtitle{Frameworks, habits, confidence}
\author{Stephen Wolff\thanks{\url{https://www.linkedin.com/in/stephen-wolff-0x200/}}}
%\date{}



\begin{frame}
\titlepage
\end{frame}



\section{Introduction}



\begin{frame}
\frametitle{Goals}

Help you\ldots{}

\begin{itemize}
  \item Recognize and respond to cyber risk
  \begin{itemize}
    \item Adopt simple, useful frameworks to manage risk
    \item Aim to reduce risk, not eliminate it
  \end{itemize}
  \item Adopt safe cyber habits
  \begin{itemize}
    \item Plan and prepare
    \item Focus on habits, not tools
  \end{itemize}
  \item Be informed and confident in your cyber decisions
  \begin{itemize}
    \item Have a reference list of trustworthy tools and information
  \end{itemize}
\end{itemize}

\end{frame}



\begin{frame}
\frametitle{Outline}

\begin{enumerate}
  \item Fundamentals
  \item Your security
  \item Toolkit overview
  \item Conclusion
\end{enumerate}

\end{frame}



\section{Fundamentals}



\begin{frame}
\frametitle{Fundamentals}

\begin{block}{Outline:}

\begin{itemize}
  \item Definitions
  \item Axioms
  \item Frameworks
\end{itemize}
\end{block}

\end{frame}



\begin{frame}
\frametitle{Definitions}

What is cybersecurity?

\begin{itemize}
  \item Identify, assess, address, and monitor relevant cyber threats
\end{itemize}

\end{frame}



\subsection{Axioms}

\begin{frame}
\frametitle{Axioms}

\begin{block}{Outline:}
\begin{enumerate}
  \item The Internet is
  \item Tradeoffs exist
  \item Cyberattacks happen
\end{enumerate}
\end{block}

\end{frame}



\begin{frame}
\frametitle{The Internet}

\begin{itemize}
  \item The Internet is public
  \begin{itemize}
    \item Anyone can use it
  \end{itemize}
  \item The Internet is open
  \begin{itemize}
    \item Anyone can see Internet traffic
  \end{itemize}
  \item The Internet is forever
  \begin{itemize}
    \item \href{https://archive.org/}{Internet Archive} (Wayback Machine)
  \end{itemize}
\end{itemize}

\end{frame}



\begin{frame}
\frametitle{Tradeoffs}

\begin{itemize}
  \item Tradeoffs exist
  \begin{itemize}
    \item Security versus cost
    \item Security versus convenience
    \item Trust versus time
  \end{itemize}
  \item Tradeoffs change
  \begin{itemize}
    \item Technologies change
    \item Threats change
    \item Values change
  \end{itemize}
  \item Tradeoffs depend on context
\end{itemize}

\end{frame}



\begin{frame}
\frametitle{Cyberattacks happen}

\begin{itemize}
  \item You can't control\ldots{}
  \begin{itemize}
    \item Whether cyberattacks happen (they will)
    \item Whether you achieve perfect safety (you can't)
  \end{itemize}
  \item You can control\ldots{}
  \begin{itemize}
    \item How informed you are
    \begin{itemize}
      \item Know what to do and where to go
    \end{itemize}
    \item How prepared you are
    \item What balance \emph{you choose} among tradeoffs
    \item What laws you support (privacy, accountability)
    \item What grace you show yourself and others
    \begin{itemize}
      \item More perfect, never perfect
    \end{itemize}
  \end{itemize}
\end{itemize}

\end{frame}



\begin{frame}
\frametitle{Your turn}

\begin{block}{Which is riskier?}<1->
\begin{enumerate}[(a)]
  \item Your bff AirDrops a file to you irl
  \item A stranger sends you a file over the Internet
\end{enumerate}
\end{block}

\begin{block}{Context:}<2->
\begin{itemize}
  \item Your bff
  \begin{itemize}
    \item Every file your bff has shared with you recently has been infected with malware
    \item The file seems unusually large
  \end{itemize}
  \item The stranger
  \begin{itemize}
    \item The stranger is (allegedly) an employee at Microsoft
    \item The file is a security update whose checksums agree with those on Microsoft's website
  \end{itemize}
\end{itemize}
\end{block}

\end{frame}



\begin{frame}
\frametitle{Frameworks}

\begin{block}{Outline:}
\begin{itemize}
  \item CIA triad
  \item Risk
  \item AAA
  \item OODA loop
\end{itemize}
\end{block}

\end{frame}



\begin{frame}
\frametitle{CIA triad}

\begin{itemize}
  \item Confidentiality
  \begin{itemize}
    \item Unauthorized entities can't access resources
    \item Opposite: disclosure
  \end{itemize}
  \item Integrity
  \begin{itemize}
    \item Resources are reliable
    \item Opposite: alteration
  \end{itemize}
  \item Availability
  \begin{itemize}
    \item Authorized entities can access resources
    \item Opposite: denial of service
  \end{itemize}
\end{itemize}

\begin{alertblock}{How to use it:}

\begin{itemize}
  \item To understand and assess security
\end{itemize}
\end{alertblock}

\end{frame}



\begin{frame}
\frametitle{Risk}

\begin{itemize}
  \item Asset
  \begin{itemize}
    \item Something of value
    \item Can be physical, digital, or intangible
  \end{itemize}
  \item Vulnerability
  \begin{itemize}
    \item A weakness
  \end{itemize}
  \item Threat
  \begin{itemize}
    \item A potential danger
  \end{itemize}
  \item Risk
  \begin{itemize}
    \item A potential loss that occurs when a threat exploits a vulnerability of an asset
  \end{itemize}
\end{itemize}

\begin{alertblock}{How to use it:}

\begin{itemize}
  \item To understand and manage risk
\end{itemize}
\end{alertblock}

\end{frame}



\begin{frame}
\frametitle{AAA}

\begin{itemize}
  \item Identification
  \begin{itemize}
    \item Who are you?
  \end{itemize}
  \item Authentication
  \begin{itemize}
    \item How do I know you are who you say you are?
  \end{itemize}
\item Authorization
  \begin{itemize}
    \item What are you allowed to do?
  \end{itemize}
  \item Accounting
  \begin{itemize}
    \item What did you do?
  \end{itemize}
\end{itemize}

\begin{alertblock}{How to use it:}
\begin{itemize}
  \item To understand interactions with service providers
\end{itemize}
\end{alertblock}

\end{frame}



\begin{frame}
\frametitle{Factors of authentication}

\begin{itemize}
  \item Something you know
  \begin{itemize}
    \item Password
  \end{itemize}
  \item Something you have
  \begin{itemize}
    \item Cell phone
  \end{itemize}
  \item Something you are
  \begin{itemize}
    \item Face, fingerprint
  \end{itemize}
\end{itemize}

\begin{alertblock}{How to use it:}
\begin{itemize}
  \item To understand MFA
  \begin{itemize}
    \item Use two or more \emph{different} factors during authentication
  \end{itemize}
\end{itemize}
\end{alertblock}

\end{frame}



\begin{frame}
\frametitle{OODA loop}

\begin{enumerate}
  \item Observe
  \item Orient
  \item Decide
  \item Act
\end{enumerate}

\begin{alertblock}{How to use it:}
\begin{itemize}
  \item To make deliberate decisions under stress
\end{itemize}
\end{alertblock}

\end{frame}



\section{Your security}



\begin{frame}
\frametitle{Your security}

\begin{block}{Outline:}
\begin{itemize}
  \item Account access
  \item Software
  \item Devices
  \item Networks
  \begin{itemize}
    \item Home router
    \item Public Wi-Fi
  \end{itemize}
  \item Relational
  \item In public
  \item Recovery
\end{itemize}
\end{block}

\end{frame}



\begin{frame}
\frametitle{Account access: Your turn}

\begin{itemize}
  \item How many passwords do you have?
  \begin{itemize}
    \item Give an order-of-magnitude estimate: 1, 10, 100, 1000,\ldots{}
  \end{itemize}
  \item How do you manage your passwords? (Select all that apply.)
  \begin{enumerate}[(a)]
    \item I use the same password or similar passwords for most or all of my accounts.
    \item I write down my most commonly used passwords on a piece of paper and keep it on or near my computer.
    \item I save them in my Web browser (for example, using the ``remember me'' feature).
    \item I write down my passwords in a notebook and look them up when I need them.
    \item I use a password manager.
  \end{enumerate}
  \item What are the tradeoffs for each method of password management listed above?
\end{itemize}

\end{frame}



\begin{frame}
\frametitle{Account access: To do}

\begin{itemize}
  \item Use a password manager
  \item Review the security settings on your accounts
  \begin{itemize}
    \item Password strength
    \item MFA enabled?
    \item Security alerts
    \item Account-recovery options
  \end{itemize}
  \item Implement MFA on high-value and high-risk accounts
\end{itemize}

\end{frame}



\begin{frame}
\frametitle{Password best practices}

\begin{itemize}
  \item Don't reuse passwords
  \item Make passwords sufficiently long
  \item Don't use security questions
  \begin{itemize}
    \item Why not? OSINT, social engineering
  \end{itemize}
\end{itemize}

\end{frame}



\begin{frame}
\frametitle{Password managers}

\begin{block}{Benefits:}
\begin{itemize}
  \item Generates and stores unique, complex passwords
  \item Integrates with Web browsers and mobile devices
  \item Auto-fills login credentials only on valid website
\end{itemize}
\end{block}

\begin{block}{Risk:}
\begin{itemize}
  \item Keys to the kingdom
  \begin{itemize}
    \item Make your master password strong
  \end{itemize}
\end{itemize}
\end{block}

\end{frame}



\begin{frame}
\frametitle{Password managers: How to start}

\begin{enumerate}
  \item Review
  \item Choose
  \item Install
  \begin{itemize}
    \item App: on PC, phone
    \item Web-browser plug-in: on PC
  \end{itemize}
  \item Enter credentials (for each account)
  \item Update passwords
  \begin{itemize}
    \item In the password manager \emph{and} on the account
    \item Use random-password generator feature
  \end{itemize}
\end{enumerate}

\end{frame}


\begin{frame}
\frametitle{Software: To do}

\begin{itemize}
  \item Install software from trusted sources only
  \item Review the security settings and permissions of applications
  \item Uninstall unused software
\end{itemize}

\end{frame}



\begin{frame}
\frametitle{Devices: To do}

\begin{itemize}
  \item Review the security settings on your digital devices
  \begin{itemize}
    \item Login settings
    \item Lock-screen timer
    \item Encryption
  \end{itemize}
  \item Align the security settings to the device
  \begin{itemize}
    \item Tradeoff: Security versus cost and convenience
    \item Considerations:
    \begin{itemize}
      \item Device mobility
      \item Data value
    \end{itemize}
  \end{itemize}
  \item Regularly apply software patches and updates
  \begin{itemize}
    \item Automate?
  \end{itemize}
\end{itemize}

\end{frame}



\begin{frame}
\frametitle{Networks}

\begin{block}{Outline:}
\begin{itemize}
  \item Be wary of public Wi-Fi
  \item Harden your home router
\end{itemize}
\end{block}

\end{frame}



\begin{frame}
\frametitle{Public Wi-Fi}

\begin{block}{Risks:}
\begin{itemize}
  \item Rogue access points
  \item Man-in-the-middle attacks
\end{itemize}
\end{block}

\begin{block}{Habits:}
\begin{itemize}
  \item Enable MFA on sensitive accounts
  \item Avoid sensitive transactions
  \item Use HTTPS, not HTTP
  \item If HTTPS is not enough, then use a VPN
\end{itemize}
\end{block}

\end{frame}



\begin{frame}
\frametitle{Home router: To do}

\begin{itemize}
  \item Change default credentials
  \item Enable strong encryption
  \begin{itemize}
    \item Best: WPA3
    \item Second best: WPA2
    \item Don't use: WEP, WPA
  \end{itemize}
  \item Enable automatic updates
  \item Disable risky features
  \begin{itemize}
    \item Remote management
    \item Wi-Fi Protected Setup (WPS)
    \item Universal Plug and Play (UPnP)
  \end{itemize}
  \item Create a guest network
  \begin{itemize}
    \item Enable a separate, strong password
  \end{itemize}
\end{itemize}

\end{frame}



\begin{frame}
\frametitle{Relational}

\begin{block}{Outline:}
\begin{itemize}
  \item Social engineering
  \item Phishing
\end{itemize}
\end{block}

\end{frame}



\begin{frame}
\frametitle{Social engineering}

\begin{itemize}
  \item An attempt to manipulate humans
\end{itemize}

\begin{block}{Commonly used tactics:}
\begin{itemize}
  \item Scarcity
  \item Urgency
  \item Intimidation
  \item Authority
  \item Trust
  \begin{itemize}
    \item Social proof
    \item Familiarity
  \end{itemize}
  \item Likeableness
  \begin{itemize}
    \item Pity
  \end{itemize}
\end{itemize}
\end{block}

\end{frame}



\begin{frame}
\frametitle{Phishing}

\begin{itemize}
  \item Social engineering over communication technology
\end{itemize}

\begin{block}{Variants:}
\begin{itemize}
  \item Phishing
  \begin{itemize}
    \item Via e-mail
  \end{itemize}
  \item Vishing
  \begin{itemize}
    \item Via voice call
  \end{itemize}
  \item Smishing
  \begin{itemize}
    \item Via SMS
  \end{itemize}
\end{itemize}
\end{block}

\end{frame}



\begin{frame}
\frametitle{How to detect phishing}

(Often, not always!)
\begin{itemize}
  \item Is unexpected
  \item Is from an unusual source (e-mail address, phone number)
  \item Uses a generic greeting
  \item Uses social-engineering tactics
\end{itemize}

\end{frame}



\begin{frame}
\frametitle{Social engineering: To do}
\begin{itemize}
  \item Listen to your Spidey sense
  \item Slow down and question
  \item Verify by using another route
  \item If it's a phish, then report and delete
\end{itemize}

\end{frame}



\begin{frame}
\frametitle{Your turn}

You receive the following e-mail from your bank:

\begin{quote}
We have detected suspicious activity on your account.

Please use the following link to verify your identity and account activity within 24 hours or your account will be frozen: \href{foo}{verify account}
\end{quote}

You use this bank account to pay multiple important bills, several of which come due in the next few days.

\begin{block}{What do you do?}
%\begin{enumerate}[(a)]
%  \item Ignore the e-mail. It's probably fake.
%  \item Click the link and provide the requested information.
%  \item Log in to your bank's website.
%  \item Give your bank a call.
%  \item Other (explain)
%\end{enumerate}
\end{block}

\end{frame}


\begin{frame}
\frametitle{In public: To do}

\begin{itemize}
  \item Configure secure device settings in advance
  \item Don't leave your devices unattended
  \item Beware shoulder surfing
  \begin{itemize}
    \item Bystander, binoculars, camera
  \end{itemize}
  \item Beware tailgating and piggybacking
\end{itemize}

\end{frame}



\begin{frame}
\frametitle{Recovery}

A cyber incident happened!
\begin{itemize}
  \item Forgotten password
  \item Account compromise
  \item Hardware failure
  \item Ransomware
\end{itemize}
Now what?
\begin{itemize}
  \item Account resets
  \item Data backups
\end{itemize}

\end{frame}



\begin{frame}
\frametitle{Account resets}

\begin{itemize}
  \item Forgotten password
  \begin{itemize}
    \item Receive e-mail with temporary code or link
    \item Enter recovery code (issued at account creation)
  \end{itemize}
  \item Lost MFA device
  \begin{itemize}
    \item Install authenticator app on second device
    \item Capture and securely store static setup key
  \end{itemize}
  \item Locked account
  \begin{itemize}
    \item Wait
    \item Request automated password reset
    \item Contact customer support
  \end{itemize}
\end{itemize}

\end{frame}



\begin{frame}
\frametitle{Data backups}

\begin{itemize}
  \item Adopt a realistic backup routine
  \item Protects against accidental file changes, hardware failure, ransomware
  \item Options:
  \begin{itemize}
    \item External hard drive
    \item Cloud storage
  \end{itemize}
  \item Considerations:
  \begin{itemize}
    \item Privacy
    \item Availability
    \item Criticality of data
    \item Urgency of access
  \end{itemize}
\end{itemize}

\end{frame}



\section{Toolkit overview}



\begin{frame}
\frametitle{Toolkit overview}

\begin{itemize}
  \item Password manager
  \item Anti-malware software
  \item Data backup
  \item Check suspicious file or URL: \href{https://www.virustotal.com/gui/}{VirusTotal}
  \item Check data breach: \href{https://haveibeenpwned.com}{HaveIBeenPwned}
  \item Report phishing: \href{https://apwg.org/reportphishing}{APWG}
\end{itemize}

\begin{alertblock}{Caution:}
Don't upload sensitive information
\end{alertblock}

\end{frame}



\section{Conclusion}



\begin{frame}
\frametitle{Conclusion: What we aimed for (goals)}

Help you\ldots{}

\begin{itemize}
  \item Recognize and respond to cyber risk
  \begin{itemize}
    \item Adopt simple, useful frameworks to manage risk
    \item Aim to reduce risk, not eliminate it
  \end{itemize}
  \item Adopt safe cyber habits
  \begin{itemize}
    \item Plan and prepare
    \item Focus on habits, not tools
  \end{itemize}
  \item Be informed and confident in your cyber decisions
  \begin{itemize}
    \item Have a reference list of trustworthy tools and information
  \end{itemize}
\end{itemize}

\end{frame}




\begin{frame}
\frametitle{Conclusion: What we learned}

\begin{itemize}
  \item Realities
  \begin{itemize}
    \item The Internet is public, open, and forever
    \item Tradeoffs exist
    \item Cyberattacks happen
  \end{itemize}
  \item Frameworks
  \begin{itemize}
    \item CIA triad
    \item Risk
    \item AAA
    \item OODA loop
  \end{itemize}
  \item Habits
  \begin{itemize}
    \item Prepare in advance
    \item Be alert for social-engineering tactics
    \begin{itemize}
      \item Slow down, question, and verify
    \end{itemize}
    \item Make friends with a few security tools
    \item Show grace to yourself and to others
  \end{itemize}
\end{itemize}

\end{frame}



\begin{frame}
\frametitle{Conclusion: What we'll do}

\begin{itemize}
  \item Today
  \begin{itemize}
    \item Share what you learned with a friend
  \end{itemize}
  \item This week
  \begin{itemize}
    \item Enable MFA on sensitive accounts
  \end{itemize}
  \item This month
  \begin{itemize}
    \item Start using a password manager
    \item Adopt regular routines to\ldots{}
    \begin{itemize}
      \item Update software
      \item Back up data
    \end{itemize}
  \end{itemize}
  \item This year
  \begin{itemize}
    \item Configure secure settings for accounts and devices
  \end{itemize}
\end{itemize}

\end{frame}



\begin{frame}
\frametitle{Thank you}

\begin{itemize}
  \item Presentation by: Stephen Wolff
  \item Presentation available at: \url{https://shw3512.github.io/teach/cysec-g30/}
  \item LinkedIn profile: \url{https://www.linkedin.com/in/stephen-wolff-0x200/}
\end{itemize}

\end{frame}



\end{document}